La question directrice de ce travail est la conservation d'une structure engendrée lorsque sa réalisation mathématique change. L'Holographie modulaire triadique produit, au moyen d'opérations arithmétiques relevées, des états retenant simultanément l'orientation, l'incidence, la frontière et la mémoire. Sa structure discrète du continu réunit les prolongements compatibles de ces états. La réalisation archimédienne exprime des coordonnées par raffinement de cylindres ; la réalisation conjuguée conserve des relations exceptionnelles d'incidence. La généalogie précède donc les valeurs et les formes sous lesquelles elles sont représentées. C'est la thèse de départ que le développement formel suivant expose avec ses opérations et ses dépendances.

Le problème géométrique propre à l’article commence lorsqu’une sortie
positive de cette construction devient une séparation entre nœuds,
un poids de réseau, une coordonnée de Plücker, un moment d’une mesure, un coefficient
d’énergie ou une direction de déformation réticulaire. Pour justifier ce passage,
il faut déterminer quelle information retrouve chaque représentation, comment
elle se compare aux autres et ce que conserve son raffinement. L’intérêt
d’une réalisation positive augmente de manière décisive lorsqu’elle admet un
inverse marqué : la géométrie cesse alors de n’offrir qu’une évaluation
de l’état et conserve désormais l’information utilisée par ses lecteurs.
La marque comprend l’ordre des positions, l’échelle totale et les données
enrichies qu’exige le lecteur particulier.

L’apport central s’organise autour de trois opérations. La première
est la reconstruction : la charge et la mémoire récupèrent les coordonnées
du registre, et les quotients de mineurs récupèrent ses séparations.
La seconde est le raffinement compatible : une partition se subdivise
en conservant ses extrémités, ses marques et une archive exacte du détail.
La troisième est le transport des structures : une forme, un opérateur ou un
lecteur est transporté à une autre représentation au moyen d’une application explicite. Ces
opérations permettent de démontrer la conservation des résidus, de l’énergie effective,
des moments et des spectres dans les domaines concrets où chaque identité a
un sens. Le qualificatif positif acquiert ainsi plusieurs réalisations
mathématiques précises, articulées par la reconstruction de la même donnée.

La géométrie positive possède ses propres antécédents, qui permettent de situer précisément le niveau de cette contribution. Les mesures de frontière des réseaux orientés paramètrent des cellules de la grassmannienne totalement non négative dans le travail de Postnikov~\cite{postnikov} ; Scott décrit la structure d'algèbre amassée de l'anneau de coordonnées grassmannien~\cite{scott}. Arkani-Hamed et Trnka introduisent l'amplituèdre dans la description géométrique des amplitudes de la théorie supersymétrique \(\mathcal N=4\) de Yang--Mills à la limite planaire~\cite{amplituhedron}. La formulation des géométries positives caractérise ensuite leurs formes canoniques par des pôles logarithmiques et des résidus récursifs de frontière~\cite{positive}. Le présent développement compose ces langages de réalisation avec des données préalablement engendrées par HMT. Son résultat distinctif est une collection marquée, récupérable et raffinable, accompagnée des applications reliant sa forme positive et ses lectures matérielles.

Le raffinement étend en outre la portée dimensionnelle. Le nombre de nœuds croît par subdivisions nonadiques et permet de construire des matrices positives de tout rang et de toute dimension externe finis compatibles avec ce nombre. La preuve repose sur les déterminants de Vandermonde et sur des formes de moments de rang plein. En rang un, on obtient la géométrie convexe complète, avec couverture par simplexes et résidus en toute dimension finie. Pour les rangs supérieurs, on construit une collection explicite d'images amplituèdrales et l'on préserve leurs drapeaux marqués. La couverture globale d'un amplituèdre de rang supérieur constitue un énoncé de portée différente du théorème d'existence et de compatibilité de cette collection ; les preuves ultérieures conservent cette distinction dans leur formulation.

La dimension physique est introduite par les opérateurs correspondants. Le calendrier tritique distingue rotation elliptique, ouverture hyperbolique et transport parabolique de mémoire. Leurs réalisations hamiltoniennes conservent la forme symplectique et permettent de calculer la transformation de Legendre ainsi que les termes de connexion issus d'une coordinatisation variable. Dans le secteur de jauge, l'échelle aire--volume, la mesure commune et les formes de réflexion précèdent la reconstruction de l'hamiltonien de Yang--Mills. L'argument spectral est présenté avec la coercivité uniforme, le témoin qui atteint le premier niveau positif et les opérateurs d'entrelacement qui le préservent. La relation avec les formes canoniques est établie au moyen de leurs lecteurs et de leurs résidus, tandis que le saut appartient à l'hamiltonien et à son espace physique. Cette hiérarchie situe chaque résultat à l'endroit où sa démonstration produit effectivement la conclusion.

L'exposé incorpore d'abord la construction formelle nécessaire pour obtenir l'état et son registre. Il développe ensuite les réalisations positives, l'extension dimensionnelle, l'énergie et la dualité réticulaire ; il réunit enfin les réalisations hamiltoniennes et le développement de jauge. Les programmes de vérification accompagnent des identités déterminées et des contrôles négatifs explicites. La démonstration symbolique conserve la portée générale ; l'exécution exacte vérifie les spécialisations déclarées et leur correspondance avec les formules imprimées.
